\chapter{A Calculus for Worst-Case Cost Analysis}
\label{chap:4}
\raggedbottom

% Chapter-local layout settings; all text and mathematical content are unchanged.
\begingroup
\renewcommand{\topfraction}{0.9}
\renewcommand{\bottomfraction}{0.8}
\renewcommand{\textfraction}{0.1}
\renewcommand{\floatpagefraction}{0.75}
\setlength{\textfloatsep}{12pt plus 2pt minus 2pt}
\setlength{\intextsep}{10pt plus 2pt minus 2pt}
\setlength{\floatsep}{10pt plus 2pt minus 2pt}
\makeatletter
\setlength{\@fptop}{0pt}
\setlength{\@fpsep}{16pt plus 2pt}
\setlength{\@fpbot}{0pt plus 1fil}
\makeatother

This chapter presents $\lambda^g_{\mathsf{worst}}$, a calculus for worst-case cost analysis of productive non-terminating programs using 
the later modality and graded monads.
Specifically, $\lambda^g_{\mathsf{worst}}$ adapts fragments from the guarded $\lambda$-calculus \citep{cloustonGuardedLambdaCalculusProgramming2017a}
and $\lambda$-amor \citep{rajaniUnifyingTypetheoryHigherorder2021a}.

Firstly, in Section~\ref{sec:cost-foundation} we examine how the interaction between grades and the later modality yields the notion of \emph{time-dependent cost}, 
the foundation of cost analysis in $\lambda^g_{\mathsf{worst}}$. 
We then present the syntax, type system, and operational semantics in
Sections~\ref{sec:worst-case-syntax}, \ref{sec:worst-case-type-system},
and~\ref{sec:worst-case-semantics}.
Finally, we show the cost analysis features of the calculus through example programs in Section~\ref{sec:worst-case-program-examples}.
\FloatBarrier

\section{Time-Dependent Cost}
\label{sec:cost-foundation}
As introduced in Section~\ref{sec:bg-guarded}, guarded recursion provides a discrete notion of \emph{logical time} through the later modality $\triangleright$. 
We can think of this as an implicit clock at the type level, where each occurrence of $\triangleright$ represents a delay of one logical time step, corresponding to a single clock tick.
Combined with the graded monad $\tmonad{I}{\tau}$, this notion lets us express costs that vary from step to step.

Firstly, we use an abstract parameter $i\in\mathbb N$ to denote the time. Advancing one time step corresponds to moving from $i$ to $i+1$. 
Our aim is to specify non-uniform costs that change with $i$. For instance, if $i$ starts at zero, a cost of $i+2$ 
gives costs 2, 3, 4 \dots as $i$ advances to $i+1$, $i+2$ \dots, while a constant cost like 3 remains unchanged as $i$ increments.

Similar to $\lambda$-amor, we use an \emph{index} language to represent quantitative information contained in different types.

\begin{defn}(Index)
\label{def:index}
$
  \\
  \begin{aligned}
    I,J ::= {}
    \kappa \mid N \mid R
    \mid I+J \mid I-J
    \mid \textstyle\sum_{\kappa<s} I
    \mid \mathsf{clk}
    \mid \mathsf{S}^s(I)
  \end{aligned}
$
\end{defn}

Indices are written as $I$ and $J$, and they represent \emph{symbolic} costs that can be instantiated by $i$. Below we define the time interpretation of indices.

\begin{defn}(Time interpretations of indices)
\label{def:time-indices}
  \begin{gather*}
    \begin{alignedat}{2}
      \clk(i) &= i
        &\qquad N(i) &= N \\
      R(i) &= R
        &\qquad (\mathsf{S}^s I)(i) &= I(i+s) \\
      (I+J)(i) &= I(i) + J(i)
        &\qquad (I-J)(i) &= I(i) - J(i)
    \end{alignedat}
    \\[0.8em]
    \begin{aligned}
       (\textstyle\sum_{\kappa<s}I)(i)
        &=(\sum_{k=0}^{s-1}I[k/\kappa])(i).
    \end{aligned}
  \end{gather*}
\end{defn}
All indices can be considered as functions from the time 
$i$ to a \emph{nonnegative} real number. The index $\clk$ represents the clock, i.e., the value of the time $i$. As the clock ticks, $i$ 
increments and the interpretation of $\clk$ changes. In addition, $i$ can be regarded as a Kripke world, where indices get their corresponding values at the current world.

The indices $N$ and $R$ represent constant functions on $i$. $N$ is a natural-number constant, and $R$ is a \emph{nonnegative} real-number constant. Consequently, their interpretations follow their original values.
The interpretations of addition and subtraction between indices are defined compositionally.
The $\mathsf{S}^{s}I$ index denotes the cost $I$ when the time is shifted forward by $s$, so the interpretation of the inner $I$ becomes $I(i+s)$.
$s$ is the shift index, and we write its interpreted value directly as $s$ instead of $s(i)$,
meaning that $s$ is a constant whose interpretation is always itself.\footnote{Throughout this thesis, we reserve the symbol $s$ for time-independent constant indices.} In addition, $s$ is restricted to be a natural number within a shift expression to
ensure that the shift operation only increments a fixed amount of time.
Intuitively, $\mathsf{S}^sI$ allows us to express costs that take place in the future.

The symbol $\kappa$ denotes an index variable whose value depends on an assignment to that variable. We therefore omit its time interpretation here.

Finally, the index $\textstyle\sum_{\kappa<s} I$ represents bounded summation that adds $I$ 
from $0$ up to $(s - 1)$. In particular, $I$ can contain the variable $\kappa$ bound by the sum, and $s$ needs to be a natural number constant. Therefore, the symbolic costs can be summed directly, 
i.e., $I[0/\kappa] + \ldots + I[(s-1)/\kappa]$ or $\sum_{k=0}^{s-1}I[k/\kappa]$.
Moreover, the sum will simply be $0$ if $s$ is zero.

We now consider how the graded monad interacts with the later modality when combined in a single type system.
For example, a computation of type $\tmonad{\clk}{\tau}$ produces a \emph{symbolic} cost $\clk$, and the real interpreted cost will be the value of 
$i$. If the time starts from 0, then the cost is 0. By contrast, a computation of type $\triangleright\tmonad{\clk}{\tau}$ means that it
will be executed one step later, and the time-interpreted cost will be $i+1 = 1$. 
Similarly, the cost will be $n$ for a computation of type $\triangleright^n \tmonad{\clk}{\tau}$, where $n$ is the number of nested later ($\triangleright$) types. 

In contrast, computations of types $\tmonad{3}{\tau}$ and $\triangleright^n \tmonad{3}{\tau}$ do not distinguish their costs.

\FloatBarrier

\section{Syntax}
\label{sec:worst-case-syntax}
In this section, we present the full syntax of $\lambda^g_{\mathsf{worst}}$, providing a general description of the surface constructs included in the calculus.
Since the calculus combines parts of $\lambda$-amor and the guarded $\lambda$-calculus, the formal definition is similar.

\subsection{Indices, Sorts and Constraints}
The index language should be familiar, as described in Section~\ref{sec:cost-foundation}. 
A sort determines the range of indices, and indices may belong to two possible sorts: natural numbers $\mathbb{N}$, and 
nonnegative real numbers $\mathbb{R}_{\ge 0}$. Furthermore, $\lambda^g_{\mathsf{worst}}$ has constraints over indices using the 
predicates $=$ and $<$. There are also conjunction $\land$ and disjunction $\lor$ between constraints. 
Since indices depend on the time $i$, predicates are determined by $i$ as well. The interpretation of predicates is defined as follows:
\begin{defn}[Interpretation of Predicates]
\label{def:constraint-entailment}
For any two indices $I$ and $J$, we define
\[
\begin{aligned}
  (I < J)
  &\quad\triangleq\quad
  \forall i \in \mathbb{N}.\; I(i) < J(i),\\
  (I = J)
  &\quad\triangleq\quad
  \forall i \in \mathbb{N}.\; I(i) = J(i).
\end{aligned}
\]
\end{defn}

\begin{figure}[!htbp]
  \centering
  \begin{align*}
    &\text{Types}
    &\tau &::=
      \alpha \mid \mathbf{1} \mid \mathbf{b}
      \mid \tau_1 \to \tau_2
      \mid \tau_1 \times \tau_2
      \mid \tau_1 + \tau_2 \\
    &&&\quad
      \mid \mu\alpha.\tau
      \mid \triangleright \tau
      \mid \tmonad{I}{\tau}
      \mid L^s \tau
      \mid \forall\kappa:S.\tau
      \\[6pt]
    &\text{Terms}
    &e &::=
      x \mid \langle\rangle \mid b
      \mid \tlam{x}{e}
      \mid nil \mid \Lambda.e
      \mid e_1 e_2 \\
    &&&\quad
      \mid \fold e \mid \unfold e
      \mid e_1 \circledast e_2
      \mid \delay e \\
    &&&\quad
      \mid \tuple{e_1}{e_2}
      \mid \mathsf{fst}\,e
      \mid \mathsf{snd}\,e
      \mid \mathsf{inl}\,e
      \mid \mathsf{inr}\,e \\
    &&&\quad
      \mid e\;[\,]
      \mid e_1::e_2
      \mid (\case{e}{x}{e_1}{y}{e_2}) \\
    &&&\quad
      \mid (\lmatch{e}{e_1}{h}{t}{e_2})
      \mid \ret e \\
      &&&\quad
      \mid \;\tick{I}
      \mid \bind{x}{e_1}{e_2}
      \mid \pull e
      \\[6pt]
    &\text{Indices}
    &I,J &::=
      \kappa \mid N \mid R \mid \clk
      \mid I+J \mid I-J \\
    &&&\quad
      \mid \textstyle\sum_{\kappa<s} I
      \mid \mathsf{S}^s(I)
      \\[6pt]
    &\text{Sorts}
    &S &::=
      \mathbb{N} \mid \mathbb{R}_{\ge 0}
      \\[6pt]
    &\text{Constraints}
    &C &::=
      I < J \mid I = J \mid C \land C \mid C \lor C
  \end{align*}
  \caption{Syntax of $\lambda^g_{\mathsf{worst}}$}
  \label{def:worst-case-syntax}
\end{figure}

From Definition~\ref{def:time-indices} and the pointwise interpretation of predicates, we can easily derive equations about the shift operation: for any index $I$, $\mathsf{S}^0 I= I$ (and we write $\mathsf{S}I$ for $\mathsf{S}^1 I$), 
and for any constant natural number $s, s'$, $\mathsf{S}^{s}(\mathsf{S}^{s'}I) = \mathsf{S}^{s+s'}I$.
In addition, the shift operation distributes over standard arithmetic expressions:
\begin{prop}[Distributivity of Shift]
\label{lem:shif-dist}
For any well-formed indices $I, J$, and $s$, if $ s$ is a natural number constant, then the following equations hold:
\begin{align*}
  &(1) \qquad \mathsf{S}^s(I+J)
  = \mathsf{S}^s(I)+\mathsf{S}^s(J)\\
  &(2)\qquad \mathsf{S}^s(I-J)
  = \mathsf{S}^s(I)-\mathsf{S}^s(J)\\
  &(3)\qquad \text{If } s' \text{ is a natural number constant}, \text{ then }\mathsf{S}^s\left(\textstyle\sum_{\kappa<s'}I\right)
  = \textstyle\sum_{\kappa<s'}\mathsf{S}^s(I)
\end{align*}
\end{prop}
\begin{proof}
  See Proposition~\ref{appa:shift-dist}.
\end{proof}

We do not fix a proof system for constraints in $\lambda^g_{\mathsf{worst}}$. Instead, we use a context $\Delta$ to collect constraint assumptions, and an abstract validity relation, written $\Delta \models C$, which states that $C$ follows from $\Delta$. 
Thus, all constraint reasoning is delegated to an oracle, which in practice can be implemented by a backend such as an SMT solver. 
The next section shows how this relation is used for type formations.

\subsection{Types}
$\lambda^g_{\mathsf{worst}}$ contains the types of the ordinary simply typed $\lambda$-calculus: unit, base, function,
product and sum types. Following $\lambda$-amor, we also include types that contain quantitative information.
The list type $L^s \tau$ carries an index $s$, representing its length. Furthermore, the length index is
restricted to be a natural number constant for correctness. The graded monad $\tmonad{I}{\tau}$, as previously mentioned, denotes computations returning data of type $\tau$, producing $I$ units of symbolic cost.
The polymorphic type $\forall\kappa\,{:}S.\tau$ allows types to be parameterised by indices, such as a cost-polymorphic monad
$\forall \kappa:\mathbb{N}.\;\tmonad{\kappa}{\tau}$.

Based on the guarded $\lambda$-calculus, $\lambda^g_{\mathsf{worst}}$ also includes the recursive type $\mu \alpha.\tau$.
We will see how this recursive type alongside the later type ($\triangleright \tau$) enables guarded recursion in Section~\ref{sec:worst-case-program-examples}.

For simplicity, we omit $\lambda$-amor's existential index quantification and quantification over types. 
In addition, $\lambda$-amor incorporates constraint types to ensure certain constraints hold for programs. 
However, we find that using constraints for type and subtype formations is enough for $\lambda^g_{\mathsf{worst}}$.

\subsection{Terms}
Most term constructs presented in Figure~\ref{def:worst-case-syntax} are standard. We therefore focus on the less conventional constructs introduced in our calculus, particularly those related to 
guarded recursion and cost analysis.
$\delay$ and $\circledast$ are the applicative functor operations of the later type ($\triangleright \tau$) \citep{mcbrideP08}.
The $\delay$ constructor introduces a term of the later type, and the $\circledast$ operator composes two terms of later type as a function application.
Lists have the constructors $nil$ and cons (::), with a pattern matching operation $\mathsf{match}$ as the eliminator.

We use iso-recursive types in $\lambda^g_{\mathsf{worst}}$, where a recursive type and its unfolding are isomorphic.
This is achieved by the $\fold$ and $\unfold$ operators, mapping recursive types and their unfoldings to each other.
Universal index quantification is introduced by the index abstraction $\Lambda.e$ and eliminated by the application $e\;[\,]$.
Note that index variables are omitted in these terms, following $\lambda$-amor, where index variables cannot occur in terms.

There are several constructors for the type $\tmonad{I}{\tau}$, including the standard unit and bind operations of monads: $\ret e$ and $\bind{x}{e_1}{e_2}$.
The tick operator ($\tick{I}$) is the atomic operation for incurring a symbolic cost of $I$.
The $\pull$ operator takes a delayed computation, namely a term of type $\triangleright \tmonad{I}{\tau}$,
executes the future cost and returns the delayed result. Crucially, $\pull$ distributes over the later modality and the graded monad, and we will 
further discuss it in the next section.

\FloatBarrier

\section{Type System}
\label{sec:worst-case-type-system}
In this section, we present the typing judgments of $\lambda^g_{\mathsf{worst}}$.
We begin with index sort typing and type well-formedness, which provide the basis for the subsequent development of subtyping and term typing.

% \begin{figure}[!htbp]
%   \centering
%   \[
%   \begin{gathered}
%     \fbox{$\Theta;\Delta \vdash I : S$} \qquad
%     \fbox{$\Psi;\Theta;\Delta \vdash \tau$} \\[0.5em] 
%     \fbox{$\Theta;\Delta \vdash \tau <: \tau'$} \qquad
%     \fbox{$\Theta;\Delta \vdash \Gamma <: \Gamma'$} \qquad
%     \fbox{$\hastype{\Theta;\Delta;\Gamma}{e}{\tau}$}
%   \end{gathered}
%   \]
%   \caption{judgment forms of $\lambda^g_{\mathsf{worst}}$}
%   \label{fig:worst-case-judgments}
% \end{figure}

\subsection{Indices and Sorts}

\begin{figure}[!htbp]
  \small
  \begin{mathpar}[\rulelayout{0.98}]
    \inferrule*[right=S-Var]
        { }
        {\Theta,\kappa:S;\Delta\vdash \kappa:S}
    \and
    \inferrule*[right=S-Nat]
        { }
        {\Theta;\Delta\vdash N:\mathbb{N}}
    \and
    \inferrule*[right=S-Real]
        { }
        {\Theta;\Delta\vdash R:\mathbb{R}_{\ge0}}
    \\
    \inferrule*[right=S-Clock]
        { }
        {\Theta;\Delta\vdash \clk: \mathbb{N}}
    \and
    \inferrule*[right=S-NatReal]
        {\Theta;\Delta\vdash I: \mathbb{N}}
        {\Theta;\Delta\vdash I: \mathbb{R}_{\ge0}}
    \\
    \inferrule*[right=S-Add]
        {\Theta;\Delta\vdash I_1:S \quad \Theta;\Delta\vdash I_2:S}
        {\Theta;\Delta\vdash I_1 + I_2:S}
    \and
    \inferrule*[right=S-Sub]
        {\Theta;\Delta\vdash I:S \quad \Theta;\Delta\vdash J:S
         \quad\Theta;\Delta\models J\leq I}
        {\Theta;\Delta\vdash I-J:S}
    \\
    \inferrule*[right=S-Sum]
        {\Theta;\Delta\vdash s:\mathbb{N} \quad
         \Theta,\kappa:\mathbb{N};\Delta, \kappa < s\vdash I:S \quad
         \Theta;\Delta\models s = \mathsf{S}(s)}
        {\Theta;\Delta\vdash\textstyle{\sum}_{\kappa<s}I:S}
    \\
    \inferrule*[right=S-Shift]
        {\Theta;\Delta\vdash I: S \quad \Theta;\Delta\vdash s:\mathbb{N}
          \quad \Theta;\Delta\models s = \mathsf{S(s)}}
        {\Theta;\Delta\vdash \mathsf{S}^s(I):S}
  \end{mathpar} 
  \caption{Sort typing rules for indices}
  \label{defn:worst-case-sort-rule}
\end{figure}

The judgment for indices is in the form $\Theta;\Delta\vdash I: S$, where $\Theta$ is the context that maps index variables $\kappa$ to their sorts $S$, and 
$\Delta$ is the context that contains constraints about indices as mentioned previously. Constraints in $\Delta$ may refer to variables in $\Theta$, so we need to use them together.

Figure~\ref{defn:worst-case-sort-rule} presents the sort rules for indices. The \textsc{S-Var, S-Nat}  and \textsc{S-Real} rules should be self-explanatory. 
The $\mathsf{clk}$ expression is treated as a natural number, representing the logical time. The \textsc{S-NatReal} rule allows a natural number index to be used as a nonnegative real number with the same value. 
For ordinary arithmetic operations, we require
both arguments to have the same sort. Since indices cannot be negative according to all the sorts ($\mathbb{N}, \mathbb{R}_{\ge0}$), the subtraction $I - J$ needs to
ensure that $J \leq I$. Here we write $J \leq I$ as the disjunction between $J < I$ and $J = I$. Similarly, in \textsc{S-Sum}, we ensure the index variable is smaller than the bound $s$ when checking 
the body expression. The constraint $s = \mathsf{S}(s)$ in \textsc{S-Sum} and \textsc{S-Shift} requires $s$ to be a time-independent constant.
Since constraints are pointwise in the time, this ensures $s$'s interpretation does not change for any $i$, which is the desired property for 
both the bounded sum and the shift index.


\subsection{Type Well-formedness}
\begin{figure}[!htbp]
  \small
  \begin{mathpar}[\rulelayout{0.99}]
    \inferrule*[right=WF-Var]
        { }
        {\Psi, \alpha;\Theta;\Delta\vdash\alpha}
    \and
    \inferrule*[right=WF-Unit]
        { }
        {\Psi;\Theta;\Delta\vdash\textbf{1}}
    \and
    \inferrule*[right=WF-Base]
        { }
        {\Psi;\Theta;\Delta\vdash \textbf{b}}
    \\
    \inferrule*[right=WF-Arrow]
        {\Psi;\Theta;\Delta\vdash\tau_1 \quad \Psi;\Theta;\Delta\vdash\tau_2 }
        {\Psi;\Theta;\Delta\vdash\tau_1\to\tau_2 }
    \and
    \inferrule*[right=WF-Prod]
        {\Psi;\Theta;\Delta\vdash\tau_1  \quad \Psi;\Theta;\Delta\vdash\tau_2 }
        {\Psi;\Theta;\Delta\vdash\tau_1\times\tau_2 }
    \\
    \inferrule*[right=WF-Sum]
        {\Psi;\Theta;\Delta\vdash\tau_1  \quad \Psi;\Theta;\Delta\vdash\tau_2 }
        {\Psi;\Theta;\Delta\vdash\tau_1+\tau_2 }
    \and
    \inferrule*[right=WF-Later]
        {\Psi;\Theta;\Delta\vdash\tau }
        {\Psi;\Theta;\Delta\vdash\triangleright\tau }
    \\
    \inferrule*[right=WF-Monad]
        {\Psi;\Theta;\Delta\vdash\tau \quad \Theta;\Delta\vdash I: \mathbb{R}_{\ge0}}
        {\Psi;\Theta;\Delta\vdash\tmonad{ I}{\tau} }
    \and
    \inferrule*[right=WF-Rec]
        {\Psi,\alpha;\Theta;\Delta\vdash\tau 
         \quad \alpha\text{ guarded in }\tau}
        {\Psi;\Theta;\Delta\vdash\mu\alpha.\tau }
    \\
    \inferrule*[right=WF-List]
        {\Psi;\Theta;\Delta\vdash\tau \quad 
        \Theta;\Delta\vdash s:\mathbb{N}
        \quad \Theta;\Delta\models s = \mathsf{S}(s)}
        {\Psi;\Theta;\Delta\vdash L^s\tau}
    \and
    \inferrule*[right=WF-Forall]
        {\Psi;\Theta,\kappa:S;\Delta\vdash\tau}
        {\Psi;\Theta;\Delta\vdash\forall\kappa:S.\tau}
  \end{mathpar} 
  \caption{Well-formedness rules for types}
  \label{def:worst-case-type-well-form}
\end{figure}
Figure~\ref{def:worst-case-type-well-form} shows the type well-formedness rules; these rules specify how types are formed correctly in $\lambda^g_{\mathsf{worst}}$.
$\Psi$ is the context that contains a set of type variables.  When checking a recursive type $\mu\alpha.\tau$, the context is extended with $\alpha$ to ensure that occurrences of this variable in $\tau$ are well-formed.
We also need the contexts $\Theta$ and $\Delta$ to check sorts and constraints since types can mention indices.
In particular, the \textsc{WF-Monad} rule requires the cost index to be a nonnegative real number. In the \textsc{WF-List} rule, we need to ensure that the 
length of a list remains fixed across time by the constraint $s = \mathsf{S}(s)$, as discussed before.

\textsc{WF-Rec} is a crucial rule in the type system. A type variable $\alpha$ is \emph{guarded} in a type $\tau$ if each
occurrence of it in $\tau$ is beneath a later modality. This brings guarded recursion, where all self-references in the recursive type need to be 
guarded.

\subsection{Subtyping}
Subtyping is an important part of our type system. In particular, it allows a cost bound to be weakened in the graded monad. 
Figure~\ref{def:worst-case-subtyping} gives the subtyping rules.

\begin{figure}[!htbp]
  \small
  \begin{mathpar}[\rulelayout{1.00}]
    \inferrule*[right=Sub-Refl]
        { }
        {\Theta;\Delta\vdash \tau <: \tau}
    \and
    \inferrule*[right=Sub-Trans]
        {\Theta;\Delta\vdash \tau_1 <: \tau_2 \quad \Theta;\Delta\vdash \tau_2 <: \tau_3}
        {\Theta;\Delta\vdash\tau_1 <: \tau_3}
    \\
    \inferrule*[right=Sub-Arrow]
        {\Theta;\Delta\vdash \tau_1' <: \tau_1 \quad \Theta;\Delta\vdash \tau_2 <: \tau_2'}
        {\Theta;\Delta\vdash \tau_1 \to \tau_2 <: \tau_1' \to \tau_2'}
    \and
    \inferrule*[right=Sub-Later]
        {\Theta;\Delta\vdash \tau <:  \tau'}
        {\Theta;\Delta\vdash \triangleright \tau <: \triangleright \tau'}
    \\
    \inferrule*[right=Sub-Product]
        {\Theta;\Delta\vdash \tau_1 <: \tau_1' \quad \Theta;\Delta\vdash \tau_2 <: \tau_2'}
        {\Theta;\Delta\vdash \tau_1 \times \tau_2 <:  \tau_1' \times \tau_2'}
    \and
    \inferrule*[right=Sub-List]
        {\Theta;\Delta\vdash \tau <: \tau'}
        {\Theta;\Delta\vdash L^s\tau <: L^s\tau'}
    \\
    \inferrule*[right=Sub-Sum]
        {\Theta;\Delta\vdash \tau_1 <: \tau_1' \quad \Theta;\Delta\vdash \tau_2 <: \tau_2'}
        {\Theta;\Delta\vdash \tau_1 + \tau_2 <:  \tau_1' + \tau_2'}
    \and
    \inferrule*[right=Sub-IndexPoly]
        {\Theta,\kappa:S;\Delta\vdash\tau<:\tau'}
        {\Theta;\Delta\vdash
          \forall\kappa:S.\tau<:\forall\kappa:S.\tau'}
    \\
    \inferrule*[right=Sub-Monad]
        {\Theta;\Delta\vdash \tau <: \tau' \quad \Theta;\Delta\models I\leq I'}
        {\Theta;\Delta\vdash \tmonad{I}{\tau} <: \tmonad{ I'}{\tau'}}
    \and
    \inferrule*[right=Sub-Rec]
        {\Theta;\Delta\vdash \tau <: \tau'
         \quad \alpha\text{ positive in }\tau,\tau'}
        {\Theta;\Delta\vdash \mu\alpha.\tau <: \mu\alpha.\tau'}
  \end{mathpar} 
  \caption{Subtyping rules for types}
  \label{def:worst-case-subtyping}
\end{figure}

The subtyping judgment is $\Theta;\Delta\vdash\tau<:\tau'$. Similar to the type well-formedness judgment, it requires $\Theta$ and $\Delta$ to ensure constraints hold for certain rules. 
We follow $\lambda$-amor and the usual treatment of subtyping: the relation $<:$ is a preorder on types that satisfies reflexivity and transitivity. These two 
properties are captured by \textsc{Sub-Refl} and \textsc{Sub-Trans}. The \textsc{Sub-Arrow} rule is standard: it is contravariant in its argument type, i.e.,
the order for arguments $\tau_1' <: \tau_1$ is the reverse of that for the whole type $\tau_1 \to \tau_2 <: \tau_1' \to \tau_2'$.

The \textsc{Sub-Monad} rule gives subtyping for the cost monad: a monad $\tmonad{I}{\tau}$ can be treated as 
$\tmonad{I'}{\tau'}$ if $\tau <: \tau'$ and $I \leq I'$. As a consequence, the cost bound of a type can be increased, which is sound because costs are upper bounds. We can safely execute a computation that incurs a lower cost than originally expected.

Subtyping recursive types is a crucial and complicated task, and we impose a restriction on it. 
To explain the \textsc{Sub-Rec} rule, we first define \emph{positive} and \emph{negative} occurrences of types, based on \cite{pierce02}.

\Needspace{7\baselineskip}
\begin{defn}
\label{def:polarity}
The \emph{positive} and \emph{negative occurrences} in a type $\tau$
are defined inductively:
\begin{enumerate}
  \item $\tau$ is a positive occurrence in $\tau$.
  \item If $\tau_1 \to \tau_2$ is a positive occurrence in $\tau$, then
        $\tau_1$ is a negative occurrence in $\tau$ and $\tau_2$ is a
        positive occurrence in $\tau$.
  \item If $\tau_1 \to \tau_2$ is a negative occurrence in $\tau$, then
        $\tau_1$ is a positive occurrence in $\tau$ and $\tau_2$ is a
        negative occurrence in $\tau$.
  \item If any of $\tau_1 \times \tau_2$, $\tau_1 + \tau_2$,
        $L^s\,\tau_1$, $\triangleright\tau_1$, $\mathbb{M}\,I\,\tau_1$,
        $\forall\kappa{:}S.\,\tau_1$ or $\mu\alpha.\tau_1$ is a
        positive occurrence in $\tau$, then each of its
        component types is a positive occurrence in $\tau$.
  \item If any of the types in (4) is a negative occurrence in $\tau$,
        then each of its component types is a negative occurrence
        in $\tau$.
\end{enumerate}
\end{defn}

We define that a type variable $\alpha$ is \emph{positive} in $\tau$ if all free occurrences of it in $\tau$ are positive, i.e., $\tau$ has no negative occurrences of $\alpha$.\footnote{We additionally require that $\alpha$ does not occur free inside any recursive type $\mu\beta.\tau'$ in $\tau$ 
where $\beta$ is not positive in $\tau'$. 
This is required since unfolding a recursive type whose bound variable occurs negatively can move $\alpha$ into a negative position. For instance,
unfolding $\mu\beta.(\triangleright\beta \to \alpha)$ gives $\triangleright(\mu\beta.(\triangleright\beta \to \alpha)) \to \alpha$, where $\alpha$ also occurs in the argument, which is a negative occurrence.}
For example, the type $\textbf{1} \to \alpha$ is a positive occurrence in itself, so $\alpha$ is a positive occurrence in it. Dually, 
$\alpha$ is a negative occurrence in the type $\alpha \to \textbf{1}$.

Due to negative occurrences, the subtyping relation does not always hold after unfolding two related recursive types. We adapt an example from \cite{zhouRevisitingIsorecursiveSubtyping2020} to our setting:
If we drop the positivity side condition from \textsc{Sub-Rec}, we can derive the following relation:
\[\mu \alpha.(\triangleright\alpha) \to \tmonad{0}{\textbf{1}} <: \mu \alpha.(\triangleright\alpha) \to \tmonad{1}{\textbf{1}} \qquad (1)\]
Unfolding the types on each side, we have
\[
  \triangleright\big(\mu \alpha.(\triangleright\alpha)\to\tmonad{0}{\textbf{1}}\big)\to\tmonad{0}{\textbf{1}} 
  <:
  \triangleright\big(\mu \alpha.(\triangleright\alpha)\to\tmonad{1}{\textbf{1}}\big)\to\tmonad{1}{\textbf{1}} 
\]
However, this relation does not hold. From \textsc{Sub-Arrow}, checking the argument types requires a reverse of (1):
\[\mu \alpha.(\triangleright\alpha) \to \tmonad{1}{\textbf{1}} <: \mu \alpha.(\triangleright\alpha) \to \tmonad{0}{\textbf{1}}\]
This is invalid since $\tmonad{1}{\textbf{1}}$ is not a subtype of $\tmonad{0}{\textbf{1}}$.

Although the preservation of subtyping is not necessary for iso-recursive types, it is a desired property needed for the meta-theory of 
$\lambda^g_{\mathsf{worst}}$ presented in the next chapter, and we will further discuss this in Section~\ref{sec:axu-lemma}.

The major limitation of restricting negative occurrences in the \textsc{Sub-Rec} rule is that it rejects all negative subtyping. However, most recursive types considered in 
$\lambda^g_{\mathsf{worst}}$ like streams only have positive occurrences. An alternative approach is adopting the Amber rules \citep{cardelli85}, which are commonly used for recursive subtyping. These rules are more permissive, allowing certain negative occurrences of recursive type variables. Nevertheless, such cases rarely arise in our type system, so we
choose the simpler restriction approach.

\subsection{Terms and Typing judgments}
The typing judgment for terms in $\lambda^g_{\mathsf{worst}}$ is written as $\hastype{\Theta;\Delta;\Gamma}{e}{\tau}$. 
$\Gamma$ is the context that maps term variables $x$ to their types, and the contexts $\Theta$ and $\Delta$ are required as usual.
Since type variables only appear in recursive types, and type substitution is used directly for unfolding, we do not need the $\Psi$ context to 
record type variables. 

\begin{figure}[!htbp]
  \small
  \begin{mathpar}[\rulelayout{1.00}]
    \inferrule*[right=T-Var]
      { }
      {\hastype{\Theta;\Delta;\Gamma, x:\tau}{x}{\tau}}
    \and
    \inferrule*[right=T-base]
      { }
      {\hastype{\Theta;\Delta;\Gamma}{b}{\textbf{b}}}
    \and
    \inferrule*[right=T-Unit]
      { }
      {\hastype{\Theta;\Delta;\Gamma}{\langle \rangle}{\textbf{1}}}
    \\
    \inferrule*[right=T-Lam]
      {\hastype{\Theta;\Delta;\Gamma, x:\tau_1}{e}{\tau_2}}
      {\hastype{\Theta;\Delta;\Gamma}{\tlam{x}{e}}{\tau_1 \to \tau_2}}
    \and
    \inferrule*[right=T-App]
      {\hastype{\Theta;\Delta;\Gamma}{e_1}{\tau_1 \to \tau_2}
       \quad \hastype{\Theta;\Delta;\Gamma}{e_2}{\tau_1}}
      {\hastype{\Theta;\Delta;\Gamma}{e_1e_2}{\tau_2}}
    \\
    \inferrule*[right=T-Pair]
      {\hastype{\Theta;\Delta;\Gamma}{e_1}{\tau_1}
       \quad \hastype{\Theta;\Delta;\Gamma}{e_2}{\tau_2}}
      {\hastype{\Theta;\Delta;\Gamma}{\tuple{e_1}{e_2}}{\tau_1\times\tau_2}}
    \\
    \inferrule*[right=T-Fst]
      {\hastype{\Theta;\Delta;\Gamma}{e}{\tau_1\times\tau_2}}
      {\hastype{\Theta;\Delta;\Gamma}{\mathsf{fst}\,e}{\tau_1}}
    \and
    \inferrule*[right=T-Snd]
      {\hastype{\Theta;\Delta;\Gamma}{e}{\tau_1\times\tau_2}}
      {\hastype{\Theta;\Delta;\Gamma}{\mathsf{snd}\,e}{\tau_2}}
    \\
    \inferrule*[right=T-Inl]
      {\hastype{\Theta;\Delta;\Gamma}{e}{\tau_1}}
      {\hastype{\Theta;\Delta;\Gamma}{\mathsf{inl}\,e}{\tau_1+\tau_2}}
    \and
    \inferrule*[right=T-Inr]
      {\hastype{\Theta;\Delta;\Gamma}{e}{\tau_2}}
      {\hastype{\Theta;\Delta;\Gamma}{\mathsf{inr}\,e}{\tau_1+\tau_2}}
    \\
    \inferrule*[right=T-Case]
      {\hastype{\Theta;\Delta;\Gamma}{e}{\tau_1+\tau_2} \quad
       \hastype{\Theta;\Delta;\Gamma,x:\tau_1}{e_1}{\tau}
       \quad \hastype{\Theta;\Delta;\Gamma,y:\tau_2}{e_2}{\tau}}
      {\hastype{\Theta;\Delta;\Gamma} {\case{e}{x}{e_1}{y}{e_2}}{\tau}}
    \\
    \inferrule*[right=T-Nil]
      { }
      {\hastype{\Theta;\Delta;\Gamma}{nil}{L^0\tau}}
    \and
    \inferrule*[right=T-Cons]
      {
      \hastype{\Theta;\Delta;\Gamma}{e_1}{\tau}
       \quad \hastype{\Theta;\Delta;\Gamma}{e_2}{L^s\tau}}
      {\hastype{\Theta;\Delta;\Gamma}{e_1::e_2}{L^{s+1}\tau}}
    \\
    \inferrule*[right=T-Match]
      {
       \hastype{\Theta;\Delta;\Gamma}{e}{L^s\tau} \\\\
       \hastype{\Theta;\Delta,s=0;\Gamma}{e_1}{\tau'} \quad
       \hastype{\Theta;\Delta,s > 0;\Gamma,h:\tau,t:L^{s-1}\tau}{e_2}{\tau'}}
      {\hastype{\Theta;\Delta;\Gamma}{\lmatch{e}{e_1}{h}{t}{e_2}}{\tau'}}
    \\
    \inferrule*[right=T-Iabs]
      {\hastype{\Theta,\kappa:S;\Delta;\Gamma}{e}{\tau}}
      {\hastype{\Theta;\Delta;\Gamma}{\Lambda.e}{\forall\kappa:S.\tau}}
    \and
    \inferrule*[right=T-Iapp]
      {\hastype{\Theta;\Delta;\Gamma}{e}{\forall\kappa:S.\tau}
       \quad \Theta;\Delta\vdash I:S}
      {\hastype{\Theta;\Delta;\Gamma}{e\;[\,]}{\tau[I/\kappa]}}
    \\
    \inferrule*[right=T-Fold]
      {\hastype{\Theta;\Delta;\Gamma}{e}{\tau[\mu\alpha.\tau/\alpha]}}
      {\hastype{\Theta;\Delta;\Gamma}{\fold e}{\mu\alpha.\tau}}
    \and
    \inferrule*[right=T-Unfold]
      {\hastype{\Theta;\Delta;\Gamma}{e}{\mu\alpha.\tau}}
      {\hastype{\Theta;\Delta;\Gamma}{\unfold e}{\tau[\mu\alpha.\tau/\alpha]}}
    \\
    \inferrule*[right=T-Next]
      {\hastype{\Theta;\Delta;\Gamma}{e}{\tau}}
      {\hastype{\Theta;\Delta;\Gamma}{\delay e}{\triangleright \tau}}
    \and
    \inferrule*[right=T-Napp]
      {\hastype{\Theta;\Delta;\Gamma}{e_1}{\triangleright (\tau_1 \to \tau_2)}  \quad
      \hastype{\Theta;\Delta;\Gamma}{e_2}{\triangleright \tau_1}}
      {\hastype{\Theta;\Delta;\Gamma}{e_1 \circledast e_2}{\triangleright \tau_2}}
    \\
    \inferrule*[right=T-Ret]
      {\hastype{\Theta;\Delta;\Gamma}{e}{\tau}}
      {\hastype{\Theta;\Delta;\Gamma}{\ret e}{\tmonad{0}{\tau}}}
    \and
    \inferrule*[right=T-Bind]
      {\hastype{\Theta;\Delta;\Gamma}{e_1}{\tmonad{ I_1}{\tau_1}}
       \quad \hastype{\Theta;\Delta;\Gamma,x:\tau_1}{e_2}{\tmonad{ I_2}{\tau_2}}}
      {\hastype{\Theta;\Delta;\Gamma}{\bind{x}{e_1}{e_2}}{\tmonad{( I_1 + I_2)}{\tau_2}}}
    \\
    \inferrule*[right=T-Tick]
      {\Theta;\Delta\vdash I:\mathbb{R}_{\ge0}}
      {\hastype{\Theta;\Delta;\Gamma}{\,\tick{ I}}{\tmonad{ I}{\textbf{1}}}}
    \and
    \inferrule*[right=T-Prepay]
      {\Theta;\Delta;\Gamma\vdash e : \triangleright \tmonad{ I}{\tau}}
      {\hastype{\Theta;\Delta;\Gamma}{\pull e}{\tmonad{\mathsf{S}( I)}{\triangleright \tau}}}
    \\
    \inferrule*[right=T-Sub]
    {\hastype{\Theta;\Delta;\Gamma}{e}{\tau}
      \quad \Theta;\Delta\vdash\tau <: \tau'}
    {\hastype{\Theta;\Delta;\Gamma}{e}{\tau'}}
    \and
    \inferrule*[right=T-Weaken]
      {\hastype{\Theta;\Delta;\Gamma}{e}{\tau}
       \quad \Theta;\Delta\vdash\Gamma'<:\Gamma}
      {\hastype{\Theta;\Delta;\Gamma'}{e}{\tau}}
  \end{mathpar}
\caption{Typing rules for terms}
\label{def:worst-case-typing-judgment}
\end{figure}

The typing rules are presented in Figure~\ref{def:worst-case-typing-judgment}. Most rules are standard, corresponding to the introduction and elimination rules for the respective type constructors.
\textsc{T-Nil} is the rule for the term $nil$ that introduces an empty list, and hence the type is $L^0 \tau$.
\textsc{T-Cons} describes the typing rule for the cons (::) term, which pushes a new element in front of a list, and the length index of the resulting list 
needs to increase by one. Consequently, pattern matching on a list has two branches, where the list is either $nil$ or cons.
The \textsc{T-Match} rule also adds more assumptions to refine existing constraints: length $s = 0$ and $s > 0$ for the corresponding branch.
\textsc{T-Fold} and \textsc{T-Unfold} are the two rules that substitute type variables bound by the recursive type.
Likewise, the term-level index instantiation $e\;[\,]$ substitutes the index $I$ for the variable $\kappa$ in the result type $\tau [I/\kappa]$.

The rules for cost analysis are more interesting. $\ret e$ turns a pure term $e$ of type $\tau$ into a monadic computation.
It neither modifies $e$ nor incurs any cost, and hence its type is $\tmonad{0}{\tau}$. 
The monadic bind sequences two computations and sums their costs, as described by the \textsc{T-Bind} rule.
\textsc{T-Tick} is the rule for the $\tick{I}$ operation, which returns a result of the unit type and produces the cost $I$.
The \textsc{T-Prepay} rule type-checks the $\pull e$ term. It distributes the later modality over the graded monad, pulling the 
delayed monad out while keeping the inner result guarded. Since we pay the future cost in advance, 
the cost becomes $\mathsf{S}(I)$, where we write $\mathsf{S}(I)$ as $\mathsf{S}^1(I)$, meaning that it is calculated by shifting the time forward once.

Finally, \textsc{T-Sub} allows a term of type $\tau$ to be assigned type $\tau'$ whenever $\tau <: \tau'$, while \textsc{T-Weaken} permits changes to the typing context. The context subtyping judgment $\Theta;\Delta\vdash \Gamma <: \Gamma'$ is defined inductively by the rules in Figure~\ref{def:worst-case-context-sub}, 
where $\cdot$ denotes the empty context. Under this relation, a typing derivation under $\Gamma'$ remains valid under $\Gamma$: the latter may contain additional variables, and each variable from $\Gamma'$ is assigned a subtype of its original type. This treatment is analogous to record subtyping.

\begin{figure}[!htbp]
  \small
  \begin{mathpar}[\rulelayout{0.90}]
    \inferrule*[right=Sub-Base]
    { }
    {\Theta;\Delta\vdash\Gamma <: \cdot}
    \and
    \inferrule*[right=Sub-Ind]
    {x:\tau \in \Gamma \quad \Theta;\Delta\vdash \tau <: \tau'
    \quad \Theta;\Delta\vdash \Gamma/x <: \Gamma'}
    {\Theta;\Delta\vdash\Gamma <: \Gamma', x:\tau'}
  \end{mathpar}
\caption{Subtyping rules for $\Gamma$}
\label{def:worst-case-context-sub}
\end{figure}

\FloatBarrier

\section{Semantics}
\label{sec:worst-case-semantics}
In this section, we present the big-step, call-by-value, fuel-based operational semantics of $\lambda^g_{\mathsf{worst}}$.
Usually, big-step semantics relates a term to its final value, but a non-terminating program such as an infinite stream has no final value.
We therefore index the semantics by a natural number $n$, called \emph{fuel}, which limits how far evaluation proceeds,
so that evaluating a term gives a finite prefix of its possibly infinite result.
Since the fuel also determines which terms are values, we first discuss the definition of values in $\lambda^g_{\mathsf{worst}}$ and connect it to the semantics.

\subsection{Values}
The judgment $e\;\mathsf{val}_n$ states that the term $e$ is a value with fuel $n$.
Since recursive definitions of terms are guarded by a $\delay$ constructor, just like recursive types are guarded by $\triangleright$,
unfolding them may produce infinitely many nested $\delay$ constructors, and anything beyond the current step lies beneath one of them.
The fuel $n$ counts how many of these nested constructors are inspected when deciding whether a term is a value.

Figure~\ref{def:values} gives the value judgments. The judgments of most constructors are defined compositionally by their components. 
For instance, \textsc{V-Pair} states that $\tuple{e_1}{e_2}$ is a value when both $e_1$ and $e_2$ are values. \textsc{V-Constant} specifies terms that are considered directly as values no matter the value of $n$.
Monads are lazy, and hence their constructors are treated as values directly.

The judgment for $\mathsf{next}\,e$ is defined by induction on $n$: with fuel $0$,
$\mathsf{next}\,e$ is a value for any term $e$, and with fuel $n+1$
it is a value when $e$ is a value with fuel $n$. Once the fuel reaches
$0$, terms inside a $\delay$ are considered suspended, and they are accepted as values without inspecting their content.

\begin{figure}[!htbp]
  \small
  \begin{mathpar}[\rulelayout{0.97}]
    \inferrule*[right=V-Inl]
      {\isval{e}{n}}
      {\isval{\mathsf{inl}\,e}{n}}
    \and
    \inferrule*[right=V-Inr]
      {\isval{e}{n}}
      {\isval{\mathsf{inr}\,e}{n}}
    \and
    \inferrule*[right=V-Pair]
      {\isval{e_1}{n} \quad \isval{e_2}{n}}
      {\isval{\tuple{e_1}{e_2}}{n}}
    \and
    \inferrule*[right=V-Fold]
      {\isval{e}{n}}
      {\isval{\fold e}{n}}
    \\
    \inferrule*[right=V-Cons]
      {\isval{e_1}{n} \quad \isval{e_2}{n}}
      {\isval{e_1::e_2}{n}}
    \and
    \inferrule*[right=V-NextBase]
      { }
      {\isval{\delay e}{0}}
    \and
    \inferrule*[right=V-NextInd]
      {\isval{e}{n}}
      {\isval{\delay e}{n+1}}
    \\
    \inferrule*[right=V-Constant]
      { e \in 
        \{\langle\rangle, b, x, \tlam{x}{e}, nil, 
          \Lambda.e, \ret e, \tick{I}, \bind{x}{e_1}{e_2},\pull e\}
      }
      {\isval{e}{n}}
  \end{mathpar}
  \caption{Value judgment}
  \label{def:values}
\end{figure}

Intuitively, if a term is a value with fuel $n$, it is also a value with less fuel.

\begin{lem}[Monotonicity of Values]
\label{lem:worst-case-value-closure}
$\\$ For any $e$, and $m, n \in \mathbb{N}$, if $\isval{e}{n}$ and $m\leq n$, then $\isval{e}{m}$.
\end{lem}
\begin{proof}
  Induction on the derivation of $\isval{e}{n}$. We present the two cases related to $\delay e$.
  % The \textsc{V-Stable} case is immediate; the cons, pair, inl, inr, and
  % fold cases follow by the IH. For \textsc{V-NextZero}, necessarily $m=0$.
  \begin{itemize}
    \item Case \textsc{V-NextBase}: since $n$ is 0, then $m = 0$, and we have the goal directly.
    \item Case \textsc{V-NextInd}: let $n=k+1$ and $e=\delay e_0$, with
    $\isval{e_0}{k}$. If $m=0$, we get the desired by \textsc{V-NextBase}. Otherwise,
    let $m=r+1$ with $r\leq k$, applying the IH with $e_0, r, k$ gives
    $\isval{e_0}{r}$. By \textsc{V-NextInd}, we have $\isval{(\delay e_0)}{m}$. \qedhere
  \end{itemize}
\end{proof}



\subsection{Evaluation}
Following $\lambda$-amor, $\lambda^g_{\mathsf{worst}}$ includes two kinds of semantics: pure evaluation for pure (non-monadic) terms and forcing evaluation for monadic computations.

\subsubsection{Pure Evaluation}
The pure evaluation judgment is defined as $e \Downarrow_n v$, evaluating a pure term into a value without incurring any cost. For convenience, we write $v$ to denote a term that satisfies the value judgment with fuel $n$.
As in the value judgment, the index $n$ is the fuel, bounding the number of nested $\delay$ constructors that evaluation may enter. 

Figure~\ref{def:worst-case-pure-eval} lists the rules of pure evaluation, which are conventional rules found in most call-by-value calculi.
Similar to the value judgment, the evaluations of $\delay e$ are defined inductively, where 
we evaluate inside the constructor recursively. If $n$ is $0$, $\delay e$ then becomes a ``thunk'' that rejects further evaluation.
Meanwhile, the rules for the $\circledast$ operator follow the same pattern, creating a thunk for function application when the fuel is 0.
Lastly, the \textsc{E-Val} rule evaluates a value with fuel $n$ to itself, as the base case for other evaluations. From Lemma~\ref{lem:worst-case-value-closure},
we can derive that a value will remain unchanged if we decrease the fuel.

For example, writing $\delay$ explicitly, the stream \texttt{ones} from Section~\ref{sec:bg-guarded} becomes:
\begin{lstlisting}[style=hs]
ones = 1 :: next ones
\end{lstlisting}
\noindent
Evaluating \texttt{ones} with fuel $0$, $1$ and $2$ gives the following values respectively:
\begin{lstlisting}[style=hs]
1 :: next ones
1 :: next (1 :: next ones)
1 :: next (1 :: next (1 :: next ones))
\end{lstlisting}
\noindent
Such recursive definitions are expressed in $\lambda^g_{\mathsf{worst}}$ by the guarded fixed-point combinator,
which we present in Section~\ref{sec:worst-case-program-examples}.

\begin{figure}[!htbp]
  \small
  \begin{mathpar}[\rulelayout{0.91}]
    \inferrule*[right=E-Val]
    {\isval{v}{n}}
    {v \Downarrow_n v}
    \and
    \inferrule*[right=E-App]
    {e_1 \Downarrow_n \tlam{x}{e} \quad e_2 \Downarrow_n v_2
     \quad e[v_2/x] \Downarrow_n v}
    {e_1e_2 \Downarrow_n v}
    \\
    \inferrule*[right=E-Pair]
    {e_1\Downarrow_n v_1 \quad e_2\Downarrow_n v_2}
    {\tuple{e_1}{e_2}\Downarrow_n\tuple{}{v_2}}
    \and
    \inferrule*[right=E-Fst]
    {e\Downarrow_n\tuple{v_1}{v_2}}
    {\mathsf{fst}\,e\Downarrow_n v_1}
    \and
    \inferrule*[right=E-Snd]
    {e\Downarrow_n\tuple{v_1}{v_2}}
    {\mathsf{snd}\,e\Downarrow_n v_2}
    \\
    \inferrule*[right=E-Inl]
    {e\Downarrow_n v}
    {\mathsf{inl}\,e\Downarrow_n\mathsf{inl}\,v}
    \and
    \inferrule*[right=E-CaseL]
    {e\Downarrow_n\mathsf{inl}\,v
      \quad e_1[v/x]\Downarrow_n v'}
    {\case{e}{x}{e_1}{y}{e_2} \Downarrow_n v'}
    \\
    \inferrule*[right=E-Inr]
    {e\Downarrow_n v}
    {\mathsf{inr}\,e\Downarrow_n\mathsf{inr}\,v}
    \and
    \inferrule*[right=E-CaseR]
    {e\Downarrow_n\mathsf{inr}\,v
      \quad e_2[v/y]\Downarrow_n v'}
    {\case{e}{x}{e_1}{y}{e_2} \Downarrow_n v'}
    \\
    \inferrule*[right=E-Cons]
    {e_1\Downarrow_n v \quad e_2\Downarrow_n l}
    {e_1::e_2\Downarrow_n v::l}
    \and
    \inferrule*[right=E-MatchNil]
    {e\Downarrow_nnil \quad e_1\Downarrow_n v}
    {\lmatch{e}{e_1}{h}{t}{e_2}\Downarrow_n v}
    \\
    \inferrule*[right=E-MatchCons]
    {e\Downarrow_n v_h::l \quad e_2[v_h/h][l/t]\Downarrow_n v}
    {\lmatch{e}{e_1}{h}{t}{e_2}\Downarrow_n v}
    \\
    \inferrule*[right=E-NextBase]
    { }
    {\delay e \Downarrow_0 \delay e}
    \and
    \inferrule*[right=E-NextInd]
    {e \Downarrow_n v}
    {\delay e \Downarrow_{n+1} \delay v}
    \\
    \inferrule*[right=E-NappBase]
    {e_1 \Downarrow_{0} \delay e_1'
     \quad e_2 \Downarrow_{0} \delay e_2'}
    {e_1 \circledast e_2 \Downarrow_{0} \delay(e_1' e_2')}
    \\
    \inferrule*[right=E-NappInd]
    {e_1 \Downarrow_{n+1} \delay (\lambda x.e)
     \quad e_2 \Downarrow_{n+1} \delay v_2
     \quad e[v_2/x] \Downarrow_n v}
    {e_1 \circledast e_2 \Downarrow_{n+1} \delay v}
    \\
    \inferrule*[right=E-Fold]
    {e \Downarrow_n v}
    {\fold e \Downarrow_n \fold v}
    \and
    \inferrule*[right=E-Unfold]
    {e \Downarrow_n \fold v}
    {\unfold e \Downarrow_n v}
    \and
    \inferrule*[right=E-Iapp]
    {e \Downarrow_n \Lambda.e' \quad e' \Downarrow_n v}
    {e\;[\,] \Downarrow_n v}
  \end{mathpar}
\caption{Pure evaluation rules}
\label{def:worst-case-pure-eval}
\end{figure}

Our notion of values and evaluation is similar to that of \cite{guattoGeneralizedModalityRecursion2018d}, in which evaluation produces a finite prefix of a potentially infinite computation. 
In Guatto's semantics, all terms evaluate to a dummy value $\mathsf{stop}$ when $n=0$, thereby discarding any unevaluated terms within a $\delay$. By contrast, our semantics preserves these suspended terms, allowing them to be evaluated further if more fuel is supplied.
To see this feature of evaluation, consider the following example:

Let $f \triangleq \lambda x.\delay(x\,\langle\rangle)$ and $\mathsf{id} \triangleq \lambda y.y$. We know $f\,\mathsf{id} = \delay(\mathsf{id}\,\langle\rangle)$ and $\mathsf{id}\,\langle\rangle = \langle\rangle$.
The evaluation of $\delay(f\,\mathsf{id})$ at $n = 1$ is: 

{
\footnotesize
\mprset{sep=1em}
\[
  \inferrule*[right=E-NextInd]
    { \inferrule*[right=E-App]
        { \inferrule*[right=E-Val]{\isval{f}{0}}{ f \Downarrow_0 f }
          \quad
          \inferrule*[right=E-Val]{\isval{\mathsf{id}}{0}}{\mathsf{id} \Downarrow_0 \mathsf{id} }
          \quad
          \inferrule*[right=E-NextBase]{ }{\delay(\mathsf{id}\,\langle\rangle)\Downarrow_0 \delay(\mathsf{id}\,\langle\rangle)} }
        { f\,\mathsf{id} \Downarrow_0 \delay(\mathsf{id}\,\langle\rangle) } }
    { \delay(f\,\mathsf{id}) \Downarrow_1 \delay(\delay(\mathsf{id}\,\langle\rangle))}
\]
}

We can give $\delay(\delay(\mathsf{id}\,\langle\rangle))$ a larger $n$ to further evaluate it:
{
\footnotesize
\mprset{sep=1em}
\[
  \inferrule*[right=E-NextInd]
    { \inferrule*[right=E-NextInd]
        { \inferrule*[right=E-App]
            { \inferrule*[right=E-Val]{\isval{\mathsf{id}}{0}}{ \mathsf{id} \Downarrow_0 \mathsf{id} }
              \and
              \inferrule*[right=E-Val]{\isval{\langle\rangle}{0}}{ \langle\rangle \Downarrow_0 \langle\rangle }
              \and
              \inferrule*[right=E-Val]{\isval{\langle\rangle}{0}}{ \langle\rangle \Downarrow_0 \langle\rangle } }
            { \mathsf{id}\,\langle\rangle \Downarrow_0 \langle\rangle } }
        { \delay(\mathsf{id}\,\langle\rangle) \Downarrow_1 \delay\,\langle\rangle } }
    { \delay(\delay(\mathsf{id}\,\langle\rangle)) \Downarrow_2 \delay(\delay\,\langle\rangle) }
\]
}

\subsubsection{Forcing Evaluation}
The forcing evaluation judgment $e \Downarrow^{c}_{n,i} v$ relates a term $e$ of type $\tmonad{I}{\tau}$
to a value $v$ of type $\tau$, and is additionally indexed by the time $i$ to interpret costs.

Figure~\ref{def:worst-case-forcing-eval} gives all the forcing evaluation rules. In particular, only constructors of monads can be forced. \textsc{E-Tick} gives the time interpretation of
$I$ in the term $\tick{I}$ and reduces it to $\langle \rangle$, hence the evaluation costs $I(i)$. The \textsc{E-Return} rule specifies that $\ret e$ just gives the pure evaluation result of 
$e$ with 0 cost. \textsc{E-Bind} composes two monadic computations $e_1$ and $e_2$ and sums their costs: it first evaluates $e_1$ purely to reach its monadic constructor, then 
forces it, and similarly for $e_2$. Of greater interest is the \textsc{E-Prepay} rule: evaluating $\pull e$ requires at least one unit of fuel to reduce inside the 
$\delay$ constructor. As we pass through a $\delay$, $i$ also increments by one locally in order to assign the future cost value.
For example, $\pull(\delay\tick{I}) \Downarrow^{I(i+1)}_{n+1, i} \delay \langle\rangle$, where the cost is $I(i+1)$. From Definition~\ref{def:time-indices}, this aligns with the typing cost $\mathsf{S}(I)$.

There is no rule for evaluating a term typed $\triangleright \tmonad{I}{\tau}$ under $n$. This keeps cost effects lazy:
if a computation is delayed, its cost should not be incurred now. To demand the data inside a delayed monadic computation, the only way is to use $\pull$ and pay for
the corresponding future cost.

\begin{figure}[!htbp]
  \small
  \begin{mathpar}[\rulelayout{0.90}]
    \inferrule*[right=E-Return]
    {e \Downarrow_n v}
    {\ret e \Downarrow_{n,i}^0 v}
    \and
    \inferrule*[right=E-Bind]
    {
      e_1 \Downarrow_n v_1 \quad 
      v_1 \Downarrow_{n,i}^{c_1}v_1'
      \quad
      e_2[v_1'/x] \Downarrow_n v_2
      \quad
      v_2 \Downarrow_{n,i}^{c_2} v_2'
    }
    {
      \bind{x}{e_1}{e_2}
      \Downarrow_{n,i}^{c_1+c_2}
      v_2'
    }
    \\
    \inferrule*[right=E-Tick]
    { }
    {\tick{ I} \Downarrow_{n,i}^{ I(i)} \langle \rangle}
    \and
    \inferrule*[right=E-Prepay]
    {
      e \Downarrow_{n+1} \delay v
      \quad
      v \Downarrow_{n,i+1}^{c}v'
    }
    {
      \pull e\Downarrow_{n+1,i}^c\delay v'
    }
  \end{mathpar}
  \caption{Forcing evaluation rules}
  \label{def:worst-case-forcing-eval}
\end{figure}

Now we can look at Theorem~\ref{thm:worst-case-soundness}, which formulates the soundness of $\lambda^g_{\mathsf{worst}}$.
The theorem states that if $e$ is a closed well-typed term that has a symbolic cost $I$ designated by its type, and if it can be evaluated purely into a value $v$, then 
the real cost $c$ of forcing $v$ to a final result $v'$ is bounded by the interpretation of $I$ at the current time 
$i$, i.e., $c \leq I(i)$.
\begin{restatable}[Soundness of $\lambda^g_{\mathsf{worst}}$]{thm}{worstsoundness}
\label{thm:worst-case-soundness}
$\\$ For any $n, i \in \mathbb{N}$, if $\,\hastype{}{e}{\tmonad{ I}{\tau}}$, $e \Downarrow_n v$ 
and $v \Downarrow^c_{n, i} v'$, then $c \leq  I(i).$ 
\end{restatable}
The soundness theorem is different from $\lambda$-amor since our cost needs to be interpreted by the logical time, 
whereas $\lambda$-amor directly uses $I$ as the bound, since its semantic cost has the same form as the cost index contained by the type.
The proof of this theorem will be discussed in the next chapter.

\FloatBarrier

\section{Example Programs}
\label{sec:worst-case-program-examples}
In this section, we present several programs in $\lambda^g_{\mathsf{worst}}$ that demonstrate the expressiveness of the 
calculus and its type-based cost analysis features.

\subsection{Guarded Recursion}
To express guarded recursive functions, a fixed point combinator is required. Here we present the guarded version of Turing's fixed point combinator 
described in \citep{cloustonGuardedLambdaCalculusProgramming2017a}.

\begin{defn}[Guarded fixed point] The fixed point is derived as follows:
\begin{align*}
  \mathsf{Rec}
    &\triangleq
      \mu\alpha.
      \bigl(\triangleright\alpha
        \to(\triangleright\tau\to\tau)\to\tau\bigr), \\
  \theta
    &\triangleq
      \tlam{y}{\tlam{f}{
        f\bigl(
          (\delay(\tlam{z}{\unfold z}))
          \circledast y
          \circledast(\delay y)
          \circledast \delay f
        \bigr)}} \\
    &: \triangleright\mathsf{Rec}
       \to(\triangleright\tau\to\tau)\to\tau, \\
  \fix &\triangleq \theta (\delay(\fold\theta))
    :(\triangleright\tau\to\tau)\to\tau.
\end{align*} 
\end{defn}

Given a function $f$, according to \textsc{E-App}, $\fix f$ reduces to
\[
          f\bigl(
            (\delay(\tlam{z}{\unfold z}))
            \circledast (\delay \fold \theta)
            \circledast(\delay \delay \fold \theta)
            \circledast \delay f
          \bigr) 
\]
From \textsc{E-NappBase} and \textsc{E-NappInd}, no matter what $n$ we are given for evaluation, it then reduces to 
\[
  f\bigl(
            \delay((\tlam{z}{\unfold z}) \,
            \fold \theta \,
            (\delay \fold \theta)\,f)
          \bigr)
\]
Suppose $n=1$. Then we can reduce inside the $\delay$. 
The inner application gives $\unfold \fold \theta$, which is then reduced to $\theta$ by \textsc{E-Unfold}.
Therefore, the whole term becomes 
\[
  f (\delay (\theta(\delay \fold \theta)f)) 
\]
This term is exactly $f (\delay \fix f)$. 
Since our call-by-value semantics evaluates arguments eagerly, different values of $n$ result in different expansions:
\begin{align*}
  &n = 0: \fix f  \\
  &n = 1: f (\delay \fix f)   \\
  &n = 2: f (\delay f (\delay\fix f)) \\
  &n = 3: f (\delay f (\delay f (\delay \fix f))) \\
  &\vdots
\end{align*}
However, when $n = 0$, $\fix f$ will reduce to the following term as the base case: 
\[  f\bigl(
            \delay((\tlam{z}{\unfold z}) \,
            \fold \theta \,
            (\delay \fold \theta)\,f)
          \bigr)\]
This is to retain the unevaluated term under $\delay$ to allow further evaluation, as mentioned before.

\subsection{Stream Processing}
We can now use the fixed point combinator to build guarded streams and derive cost analysis for 
ordinary stream processing functions.

Firstly, we define the guarded stream type and its standard functions:
\[
  \tstr{\tau}
  \triangleq
  \mu\alpha.(\tau\times\triangleright\alpha).
\]
\begin{align*}
  (::) &\triangleq \tlam{x}{\tlam{s}{\fold\tuple{x}{s}}}
  &&: \tau\to\triangleright\tstr{\tau}\to\tstr{\tau}, \\
  \head &\triangleq \tlam{s}{\mathsf{fst}\,(\unfold s)}
  &&: \tstr{\tau}\to\tau, \\
  \tail &\triangleq \tlam{s}{\mathsf{snd}\,(\unfold s)}
   &&: \tstr{\tau}\to\triangleright\tstr{\tau}.
\end{align*}
Here we overload the syntax, using the same notation (::) as list cons to construct streams.

Based on the stream constructors, we can define a family of functions that enable us to observe deeper into a stream:
\begin{flalign*}
  &\quad\mathsf{nth}_n : \tstr{\tau} \to \triangleright^n\tau && \\
  &\quad\mathsf{nth}_0 \triangleq \tlam{str}{\head str} && \\
  &\quad\mathsf{nth}_{n+1} \triangleq \tlam{str}{(\delay \mathsf{nth}_n)\circledast(\tail str)} &&
\end{flalign*}
The functions are defined inductively for convenience instead of listing all of their definitions, and $n$ 
is hardcoded as the order number of a function. In particular, 
the base case just returns the current head, and the inductive case calls previous functions on the stream tail to demand 
the next element. Since subsequent elements are guarded, the number of ``later''s in the type signature follows the exact position of the element in the stream.

We will focus on the $\tstr{\tmonad{I}{\tau}}$ type for our later examples.
It denotes a stream of computations, where each element takes $I(i)$ cost to evaluate.
For example, we can construct a stream whose cost varies for each element:
\begin{flalign*}
  &\quad\mathsf{tickTime} : \tstr{\tmonad{\clk}{\mathbf{1}}} && \\
  &\quad\mathsf{tickTime} \triangleq \fix\lambda x.\, \tick{\clk} :: x &&
\end{flalign*}
For each element, this stream simply ticks the cost representing the time.
Table~\ref{fig:ticktime-costs} shows the cost of each element. 
Intuitively, each cost aligns with the time and its element position, which can be thought of as modelling the cost 
of traversing the stream from the beginning up to the desired position.

Operationally, we can define a family of functions that allow us to evaluate the $ n$-th element of a stream, which is similar to the $\mathsf{nth}$ function:
\begin{flalign*}
&\quad\mathsf{prepayNth}_n : \forall I.\,\tstr{(\tmonad{I}{\tau})} \to \tmonad{\mathsf{S}^n(I)}{(\triangleright^n\tau)}, &&\\
&\quad\mathsf{prepayNth}_0 \triangleq \Lambda.\tlam{str}{\head str}, &&\\
&\quad\mathsf{prepayNth}_{n+1} \triangleq \Lambda.\tlam{str}{\pull((\delay (\mathsf{prepayNth}_n\,[\,]))\circledast(\tail str))} &&
\end{flalign*}
The symbolic cost $\mathsf{S}^n(I)$ records the amount of increment needed for the current time 
to obtain the delayed cost. Moreover, the index $I$ is polymorphic and can take streams of arbitrary cost.
However, there is no polymorphism for the number of nested ``laters''($\triangleright^n\tau$) in our type system, and hence we 
choose the index that is equal to $n$ for $\mathsf{S}^n(I)$.

\begin{table}[!htbp]
  \renewcommand{\arraystretch}{1.35}
  \begin{tabular}{c@{\qquad}l@{\qquad}c}
    \hline
     Position
      & Type 
      & Time Interpretation of costs \\
    \hline
    $0$
      & $\tmonad{\mathsf{clk}}{\mathbf{1}}$
      & $0$ \\
    $1$
      & $\triangleright\tmonad{\mathsf{clk}}{\mathbf{1}}$
      & $1$ \\
    $2$
      & $\triangleright\triangleright\tmonad{\mathsf{clk}}{\mathbf{1}}$
      & $2$ \\
    $3$
      & $\triangleright\triangleright\triangleright
         \tmonad{\mathsf{clk}}{\mathbf{1}}$
      & $3$ \\
    $\vdots$
      & $\qquad\vdots$
      & $\vdots$ \\
    $n$
      & $\triangleright^{n}\,\tmonad{\mathsf{clk}}{\mathbf{1}}$
      & $n$ \\
    \hline
  \end{tabular}
  \caption{Costs of elements in $\mathsf{tickTime}$, starting
    at time $i = 0$}
  \label{fig:ticktime-costs}
\end{table}

Ordinary element-wise stream processing functions can be defined easily:
\begin{flalign*}
&\quad\mathsf{mapC}: \forall I_1, I_2.\,(\tau \to \tmonad{I_1}{\tau}) \to \tstr{\tmonad{I_2}{\tau}} \to \tstr{\tmonad{(I_1 + I_2)}{\tau}} &&\\
&\quad\mathsf{mapC} \triangleq \Lambda.\Lambda.\fix\lambda m.\lambda f.\lambda str.&&\\
&\qquad(\bind{x}{\head str}{f\; x}) :: (m \circledast \delay f \circledast\tail str) &&
\end{flalign*}
$\mathsf{mapC}$ is the cost-aware version of the common stream mapping function.
The function $f$ has type $\tau \to \tmonad{I_1}{\tau}$, which is the standard type of encoding 
a computation that does some work to process its input. For streams, we need to first force the head computation
and feed its result into $f$, and the two separate costs are added by $\mathsf{bind}$.

Likewise, we can define a $\mathsf{zipWithC}$ as follows:
\begin{flalign*}
&\quad\mathsf{zipWithC}: \forall I_1, I_2, I_3.\;(\tau_1\to\tau_2\to\tmonad{I_3}{\tau_3}) &&\\
&\qquad\to\tstr{\tmonad{I_1}{\tau_1}}\to\tstr{\tmonad{I_2}{\tau_2}}
    \to\tstr{\tmonad{(I_1+I_2+I_3)}{\tau_3}} &&\\
&\quad\mathsf{zipWithC}\triangleq\Lambda.\Lambda.\Lambda.\fix\lambda z.
  \lambda f.\lambda str_1.\lambda str_2. &&\\
&\qquad E_1::(z\circledast \delay f \circledast\tail str_1\circledast\tail str_2) &&\\
&\qquad\textbf{where }E_1\triangleq
  \bind{h_1}{\head str_1}{\bind{h_2}{\head str_2}{f\,h_1\,h_2}}. &&
\end{flalign*}

The function $f$ now requires one element of each input stream, and the final cost accounts for three computations: evaluating the two heads 
and processing their values with $f$.

\subsection{Primitive Recursion}
Functions like $\mathsf{mapC}$ and $\mathsf{zipWithC}$ are standard examples of non-terminating recursive computations since they continue to perform element-wise operations on streams.
On the other hand, we can also express (primitive) recursive computations that eventually terminate, with the 
help of indexed types.

Firstly, we can derive a size-indexed natural number type using lists: 
\[\mathsf{Nat}(s) \triangleq L^s\textbf{1}\]
The natural number type is just a list of units, where the data stored is irrelevant. 

Next, we define the $\mathsf{Colist}$ type, which represents a possibly infinite guarded list:
\[
  \colist{\tau} 
  \triangleq
  \mu\alpha.\textbf{1} + (\tau\times\triangleright\alpha).
\]
\begin{align*}
  (::) &\triangleq \tlam{x}{\tlam{s}{\fold(\mathsf{inr}(\tuple{x}{s}))}}
  &&: \tau\to\triangleright\colist{\tau}\to\colist{\tau}, \\
  (\mathsf{null}) &\triangleq \fold(\mathsf{inl}\,\langle\rangle)
  &&: \colist{\tau}
\end{align*}
A colist is almost the same as a guarded stream, except that it can terminate with a unit value, represented by the ``$\mathsf{null}$'' constructor.
Similar to streams, we also use the notation (::) to denote the cons constructor.

After constructing the two types, we can use them to define functions that consume a fixed-length prefix of a stream.
For example:
\begin{flalign*}
&\quad\mathsf{takeC}:\forall s, I.\;\mathsf{Nat}(s) \to \tstr{\tmonad{I}{\tau}} 
  \to \tmonad{(\textstyle\sum_{\kappa<s}\mathsf{S}^\kappa I)}{\colist{\tau}} &&\\
&\quad\mathsf{takeC} \triangleq \fix \lambda f.\Lambda.\Lambda.\lambda n.\lambda str. &&\\
&\qquad\mathsf{match}\;n\;\mathsf{with} &&\\
&\qquad\mid\;nil \mapsto \ret (\mathsf{null}) &&\\
&\qquad\mid\;\_::n'\mapsto \bind{t}{\pull(\delay(\lambda g. g\;[\,]\;[\,]) \circledast f \circledast\tail str\circledast\delay n')}{} &&\\
&\qquad\qquad\qquad\qquad\bind{h}{\head str}{}\ret(h::t) &&
\end{flalign*}

The $\mathsf{takeC}$ function copies the first $s$ elements of the input stream into a colist. It forces each element and stores the results,
and hence the final type is a monad that sums all the costs and returns a pure colist. 

$\mathsf{Nat}(s)$ is indeed a list, and we can pattern match on 
it to achieve the effect of induction on a natural number. In the base ($nil$) case, the function returns $\mathsf{null}$ to end the colist. $\mathsf{ret}$ gives cost 0, which is consistent with the type signature since 
$\textstyle\sum_{\kappa<0}\mathsf{S}^\kappa I = 0$. In the recursive case, we use $\pull$ to access the recursive call that is originally guarded, 
bind its result and the head computation to construct the final colist. Note that the recursive call is also polymorphic, and we need to use an additional function $\delay(\lambda g. g\;[\,]\;[\,]$) to instantiate it.
Consequently, the natural number's size is instantiated with $s-1$, and the cost of the $\pull$ is $\mathsf{S}(\textstyle{\sum}_{\kappa<(s-1)}\mathsf{S}^{\kappa} I)$. 
Together with the cost of binding the stream head, the final cost is $\mathsf{S}(\textstyle{\sum}_{\kappa<(s-1)}\mathsf{S}^\kappa I) + I$.

Since we know $s$ is a time-independent natural number, we can rewrite the cost to match the one presented in the type signature:
\[
\begin{aligned}
\mathsf{S}\Bigl(\sum_{\kappa<s-1}\mathsf{S}^\kappa I\Bigr) + I 
  &= \Bigl(\sum_{\kappa<s-1}\mathsf{S}^{\kappa+1} I\Bigr) + I && (\text{by Proposition~\ref{lem:shif-dist} and } \mathsf{S}(\mathsf{S}^{\kappa}) = \mathsf{S}^{\kappa+1})\\
  &= \Bigl(\sum_{\kappa<s-1}\mathsf{S}^{\kappa+1} I\Bigr) + \mathsf{S}^0 I && (I = \mathsf{S}^{0}I) \\
  &= \sum_{\kappa<s}\mathsf{S}^\kappa I && (\text{reindex } s-1 \text{ and } \kappa +1) 
\end{aligned}
\]

Overall, indexed types such as $\mathrm{Nat}(s)$ let us process a finite prefix
of an infinite type like the guarded stream, and the total cost is a finite sum. Since
a stream never ends, the results are collected in a data structure that can terminate,
such as a colist, which also agrees with the structure of guarded recursive types.

\clearpage
\endgroup
